% Part: incompleteness
% Chapter: theories-computability
% Section: computably-axiomatizable

\documentclass[../../../include/open-logic-section]{subfiles}

\begin{document}

\olfileid{inc}{tcp}{cax}

\olsection{\printtoken{S}{axiomatizable} उपपत्ती}

उपपत्ती $\Th{T}$ साठी स्वयंसिद्धकांचा संगणनक्षम संच~$A$ असेल,
तर तिला \emph{!!{axiomatizable}} म्हणतात.
($A$ हा $\Th{T}$ चा स्वयंसिद्धकसंच आहे याचा अर्थ
$T = \Setabs{!A}{A \Proves !A}$.) नैसर्गिक संख्यांचे कोणतेही
``युक्तिसंगत'' स्वयंसिद्धकीकरण हा गुणधर्म पूर्ण करेल.
विशेषतः, सांत स्वयंसिद्धकसंच असलेली कोणतीही उपपत्ती
!!{axiomatizable} असते.

\begin{lem}
$\Th{T}$ ही !!{axiomatizable} असेल, तर $\Th{T}$ ही
!!{computably
enumerable} असते.
\end{lem}

\begin{proof}
$A$ हा $\Th{T}$ चा संगणनक्षम स्वयंसिद्धकसंच आहे असे समजू.
$!A \in T$ असल्याचे होकारार्थी प्रमाण शोधण्यासाठी
स्वयंसिद्धकांपासून $!A$ ची !!a{derivation} शोधत राहा;
ते वाक्य सदस्य नसेल, तर हा शोध थांबेलच असे नाही.

थोडे वेगळ्या शब्दांत, $!A$ हे $\Th{T}$ मध्ये असते तेव्हाच
आणि तरच $A$ मधील $!B_1$, \dots, $!B_k$ अशा स्वयंसिद्धकांची
सांत यादी आणि प्रथम-क्रम तर्कशास्त्रात
$(!B_1 \land \dots \land !B_k) \lif !A$ ची
!!a{derivation} असते. पण ``असा काही \dots आहे की \dots''
या रूपातील व्याख्या असलेला कोणताही संच !!{c.e.} असतो,
जर दुसऱ्या ``\dots'' मधील अट संगणनक्षम असेल, हे आपण आधीच जाणतो.
\end{proof}

\end{document}
